% GENERATED FILE. Edit canonical lessons, then run npm run generate:books.
% canonical-source-hash: fnv1a64:b7f9fae1b9c48b10
% canonical-lessons: UR-C524-fehrist, UR-C524-musavvada, UR-C524-chhapvana, UR-C524-tadbir, UR-C524-banvana

\chapter{Making a Plan}
\label{ch:ur-making-a-plan}
\hlchaptermodality{524}

\begin{chapteropening}
\textbf{By the end of this chapter:} \emph{I can talk about plans and what happens if they change.}

Complete the last lesson of chapter 524: I can talk about plans and what happens if they change.
\end{chapteropening}

\section[\texorpdfstring{\ur{فہرست} (fehrist)}{فہرست (fehrist)}]{\ur{فہرست} (fehrist) — a list}

\label{lesson:UR-C524-fehrist}



\begin{quote}
Before the new one: say the Urdu for a cheer, then the Urdu for to leap over.
\end{quote}

\subsection*{You'll want to know: \ur{فہرست}}
\textbf{\ur{فہرست}} — \emph{fehrist} — \textquotedblleft{}a list\textquotedblright{}.

\begin{quote}
\ur{اگر} \ur{فہرست} \ur{تیار} \ur{ہے} \ur{تو} \ur{مجھے} \ur{دیجیے۔}
\end{quote}

\begin{quote}
\emph{agar fehrist tayyār hai to mujhe dījiye.}
\end{quote}

\begin{quote}
If the list is ready, please give it to me.
\end{quote}

Feminine.

\begin{cousinweb}
\textbf{\ur{فہرست}} (\emph{fehrist}) is usually traced to Persian \emph{fihrist}, \textquotedblleft{}a list, an index\textquotedblright{}.
\end{cousinweb}

\begin{grammarlens}[title={What you've built}]
One more word for saying what happens if.
\end{grammarlens}

\subsection*{Your turn}
\begin{itemize}
\raggedright
  \item \emph{Say it:} \emph{fehrist}
  \item \emph{Say it:} \emph{fehrist}, once more
  \item \emph{Say it:} \emph{agar fehrist tayyār hai to mujhe dījiye.} — If the list is ready, please give it to me.
  \item \emph{Recall:} say the Urdu for a cheer, then the Urdu for to leap over, then say \emph{fehrist} again
\end{itemize}

\begin{tcolorbox}[breakable,colback=teal!4,colframe=teal!35!black,title={Before you move on}]
What does \ur{فہرست} mean? (\textquotedblleft{}A list\textquotedblright{}.) Say it once more.
\end{tcolorbox}
\section[\texorpdfstring{\ur{مسودہ} (musavvadā)}{مسودہ (musavvadā)}]{\ur{مسودہ} (musavvadā) — a draft (of a text)}

\label{lesson:UR-C524-musavvada}



\begin{quote}
Before the new one: say the Urdu for to leap over, then the Urdu for a list.
\end{quote}

\subsection*{You'll want to know: \ur{مسودہ}}
\textbf{\ur{مسودہ}} — \emph{musavvadā} — \textquotedblleft{}a draft (of a text)\textquotedblright{}.

\begin{quote}
\ur{اگر} \ur{مسودہ} \ur{ٹھیک} \ur{ہے} \ur{تو} \ur{ہم} \ur{اسے} \ur{چھاپیں} \ur{گے۔}
\end{quote}

\begin{quote}
\emph{agar musavvadā ṭhīk hai to ham use chhāpeṅ ge.}
\end{quote}

\begin{quote}
If the draft is all right, we will print it.
\end{quote}

Masculine.

\begin{cousinweb}
\textbf{\ur{مسودہ}} (\emph{musavvadā}) is usually traced, through Persian, to Arabic \emph{musawwada}, \textquotedblleft{}a rough copy\textquotedblright{}.
\end{cousinweb}

\begin{grammarlens}[title={What you've built}]
One more word for saying what happens if.
\end{grammarlens}

\subsection*{Your turn}
\begin{itemize}
\raggedright
  \item \emph{Say it:} \emph{musavvadā}
  \item \emph{Say it:} \emph{musavvadā}, once more
  \item \emph{Say it:} \emph{agar musavvadā ṭhīk hai to ham use chhāpeṅ ge.} — If the draft is all right, we will print it.
  \item \emph{Recall:} say the Urdu for to leap over, then the Urdu for a list, then say \emph{musavvadā} again
\end{itemize}

\begin{tcolorbox}[breakable,colback=teal!4,colframe=teal!35!black,title={Before you move on}]
What does \ur{مسودہ} mean? (\textquotedblleft{}A draft (of a text)\textquotedblright{}.) Say it once more.
\end{tcolorbox}
\section[\texorpdfstring{\ur{چھپوانا} (chhapvānā)}{چھپوانا (chhapvānā)}]{\ur{چھپوانا} (chhapvānā) — to have (something) printed}

\label{lesson:UR-C524-chhapvana}



\begin{quote}
Before the new one: say the Urdu for a list, then the Urdu for a draft.
\end{quote}

\subsection*{You'll want to know: \ur{چھپوانا}}
\textbf{\ur{چھپوانا}} — \emph{chhapvānā} — \textquotedblleft{}to have (something) printed\textquotedblright{}.

\begin{quote}
\ur{اگر} \ur{کتاب} \ur{تیار} \ur{ہوئی} \ur{تو} \ur{ہم} \ur{اسے} \ur{چھپوائیں} \ur{گے۔}
\end{quote}

\begin{quote}
\emph{agar kitāb tayyār huī to ham use chhapvāeṅ ge.}
\end{quote}

\begin{quote}
If the book is ready, we'll have it printed.
\end{quote}

Past, with \ur{نے} (he or she …): \ur{اس} \ur{نے} \ur{اپنی} \ur{کتاب} \ur{چھپوائی}, \emph{us ne apnī kitāb chhapvāī}.

\begin{cousinweb}
\textbf{\ur{چھپوانا}} (\emph{chhapvānā}) is the causative of \emph{chhāpnā}, \textquotedblleft{}to print\textquotedblright{}, which CDIAL lists under \emph{*chapp}, \textquotedblleft{}press\textquotedblright{}, a form it reconstructs (CDIAL 4994).
\end{cousinweb}

\begin{grammarlens}[title={What you've built}]
One more word for saying what happens if.
\end{grammarlens}

\subsection*{Your turn}
\begin{itemize}
\raggedright
  \item \emph{Say it:} \emph{chhapvānā}
  \item \emph{Say it:} \emph{chhapvānā}, once more
  \item \emph{Say it:} \emph{agar kitāb tayyār huī to ham use chhapvāeṅ ge.} — If the book is ready, we'll have it printed.
  \item \emph{Recall:} say the Urdu for a list, then the Urdu for a draft, then say \emph{chhapvānā} again
\end{itemize}

\begin{tcolorbox}[breakable,colback=teal!4,colframe=teal!35!black,title={Before you move on}]
What does \ur{چھپوانا} mean? (\textquotedblleft{}To have (something) printed\textquotedblright{}.) Say it once more.
\end{tcolorbox}
\section[\texorpdfstring{\ur{تدبیر} (tadbīr)}{تدبیر (tadbīr)}]{\ur{تدبیر} (tadbīr) — a way out (of a difficulty)}

\label{lesson:UR-C524-tadbir}



\begin{quote}
Before the new one: say the Urdu for a draft, then the Urdu for to have (something) printed.
\end{quote}

\subsection*{You'll want to know: \ur{تدبیر}}
\textbf{\ur{تدبیر}} — \emph{tadbīr} — \textquotedblleft{}a way out (of a difficulty)\textquotedblright{}.

\begin{quote}
\ur{اگر} \ur{کوئی} \ur{تدبیر} \ur{نہیں} \ur{ملی} \ur{تو} \ur{ہم} \ur{رک} \ur{جائیں} \ur{گے۔}
\end{quote}

\begin{quote}
\emph{agar koī tadbīr nahīṅ milī to ham ruk jāeṅ ge.}
\end{quote}

\begin{quote}
If we find no way out, we will stop.
\end{quote}

Feminine.

\begin{cousinweb}
\textbf{\ur{تدبیر}} (\emph{tadbīr}) is usually traced, through Persian, to Arabic \emph{tadbīr}, \textquotedblleft{}arranging, managing\textquotedblright{}.
\end{cousinweb}

\begin{grammarlens}[title={What you've built}]
One more word for saying what happens if.
\end{grammarlens}

\subsection*{Your turn}
\begin{itemize}
\raggedright
  \item \emph{Say it:} \emph{tadbīr}
  \item \emph{Say it:} \emph{tadbīr}, once more
  \item \emph{Say it:} \emph{agar koī tadbīr nahīṅ milī to ham ruk jāeṅ ge.} — If we find no way out, we will stop.
  \item \emph{Recall:} say the Urdu for a draft, then the Urdu for to have (something) printed, then say \emph{tadbīr} again
\end{itemize}

\begin{tcolorbox}[breakable,colback=teal!4,colframe=teal!35!black,title={Before you move on}]
What does \ur{تدبیر} mean? (\textquotedblleft{}A way out (of a difficulty)\textquotedblright{}.) Say it once more.
\end{tcolorbox}
\section[\texorpdfstring{\ur{بنوانا} (banvānā)}{بنوانا (banvānā)}]{\ur{بنوانا} (banvānā) — to have (something) made or built}

\label{lesson:UR-C524-banvana}



\begin{quote}
Before the new one: say the Urdu for to have (something) printed, then the Urdu for a way out.
\end{quote}

\subsection*{You'll want to know: \ur{بنوانا}}
\textbf{\ur{بنوانا}} — \emph{banvānā} — \textquotedblleft{}to have (something) made or built\textquotedblright{}.

\begin{quote}
\ur{اگر} \ur{نقشہ} \ur{ٹھیک} \ur{ہے} \ur{تو} \ur{ہم} \ur{گھر} \ur{بنوائیں} \ur{گے۔}
\end{quote}

\begin{quote}
\emph{agar naqshā ṭhīk hai to ham ghar banvāeṅ ge.}
\end{quote}

\begin{quote}
If the plan is right, we'll have the house built.
\end{quote}

Past, with \ur{نے} (he or she …): \ur{اس} \ur{نے} \ur{نیا} \ur{گھر} \ur{بنوایا}, \emph{us ne nayā ghar banvāyā}.

\begin{cousinweb}
\textbf{\ur{بنوانا}} (\emph{banvānā}) is the second causative of \emph{bannā}, \textquotedblleft{}to be made\textquotedblright{}, which CDIAL lists under Sanskrit \emph{vanati}, \textquotedblleft{}desires, makes ready\textquotedblright{} (CDIAL 11260): to have someone make a thing.
\end{cousinweb}

\begin{grammarlens}[title={What you've built}]
One more word for saying what happens if.
\end{grammarlens}

\subsection*{Your turn}
\begin{itemize}
\raggedright
  \item \emph{Say it:} \emph{banvānā}
  \item \emph{Say it:} \emph{banvānā}, once more
  \item \emph{Say it:} \emph{agar naqshā ṭhīk hai to ham ghar banvāeṅ ge.} — If the plan is right, we'll have the house built.
  \item \emph{Recall:} say the Urdu for to have (something) printed, then the Urdu for a way out, then say \emph{banvānā} again
\end{itemize}

\begin{tcolorbox}[breakable,colback=teal!4,colframe=teal!35!black,title={Before you move on}]
What does \ur{بنوانا} mean? (\textquotedblleft{}To have (something) made or built\textquotedblright{}.) Say it once more.
\end{tcolorbox}
